\chapter{Introduction}

\section{Kinetoplast DNA and Dyskinetoplasty}

Kinetoplast DNA (kDNA) is a defining feature of kinetoplastid protozoans, including the medically important genus \textit{Trypanosoma}. Yet some species have partially or completely lost kDNA, raising the question: what prevents loss in \teq{}?

\section{Research Question}

Why does \teq{} retain kinetoplast DNA fragments despite being dyskinetoplastic, while the closely related \tev{} has lost kDNA entirely?

\section{Competing Models}

We test three mechanistic models:

\subsection{Model A: Genetic Constraint}
kDNA retention is enforced by the absence of compensatory mutations (e.g., A273P in ATPase γ).

\textbf{Prediction:} No A273P; $\omega < 1$ (purifying selection)

\subsection{Model B: Relaxed Selection}
kDNA is retained due to weak selective pressure.

\textbf{Prediction:} No A273P; $\omega \approx 1$ (neutral evolution)

\subsection{Model C: Functional Threshold}
Nuclear compensation exists but cannot fully substitute kDNA function.

\textbf{Prediction:} Transcriptional response is incomplete

\section{Project Overview}

This computational study integrates genomic, phylogenetic, and transcriptomic evidence to distinguish these models using public data.

